Quiero expresar mi más sincero agradecimiento, en primer lugar, a mi familia. A mis padres y a mis cuatro hermanos, por su apoyo incondicional, su paciencia infinita y su constante motivación a lo largo de cada una de las etapas de este camino universitario. Su cariño, esfuerzo y confianza han sido el pilar fundamental para superar cada reto y alcanzar esta meta.

A los profesores y al conjunto de la comunidad académica de la Escuela Politécnica Superior de la Universidad Autónoma de Madrid, por la sólida formación técnica, científica y humana brindada durante estos años de carrera, que ha constituido la base indispensable para mi crecimiento formativo y profesional.

Asimismo, quiero agradecer a Xavier Alamán por su orientación académica, dedicación y valiosas observaciones durante el desarrollo y revisión metodológica de este trabajo.

Y de manera muy especial, quiero dar las gracias a mis amigos y a Amalia, por compartir conmigo tanto los momentos de esfuerzo y estudio como las alegrías y celebraciones de estos años. Por estar siempre presentes cuando ha hecho falta escuchar, desahogarse, animar y compartir la vida, haciendo que esta etapa universitaria sea un recuerdo imborrable.